\documentclass[11pt]{article}
\usepackage{amsmath,amssymb,amsthm}
\usepackage[margin=1in]{geometry}
\newtheorem{proposition}{Proposition}
\title{P6 v2: a genuinely different attempt at $\beta_2$'s arithmetic nature}
\author{}
\date{2026-09-26}
\begin{document}
\maketitle

\noindent\textbf{Scope and relation to P6\_note.tex.} The prior note
(\texttt{code/zeta\_probe/rootunity\_zeros/general/P6/P6\_note.tex}) documents that five
independent rational-approximation constructions to the Hahn--Exton $q$-Bessel value
$S(q^\ast)=0$ all land at the same decay/denominator-growth ratio $\rho=1/2$, one factor of two
below the $\rho\ge1$ threshold needed by classical $q$-irrationality criteria
(Tschakaloff/Bézivin--Bundschuh--Väänänen-type). This note does \emph{not} try a sixth
reparametrization of that same Diophantine-approximation route. Instead it asks the two
specific questions posed for this task: (i) is $q^\ast$ (rather than $\beta_2$) tractable via a
\emph{different} transcendence framework entirely (periods, Nesterenko, Schneider--Lang,
modular functions), bypassing $q$-irrationality exponents altogether; and (ii) does the bare
numerical value $q^\ast\approx0.4495$ match any known special constant, which would reveal
hidden structure. Labels PROVED / VERIFIED / HEURISTIC / UNVERIFIED-CITATION follow the same
convention as the prior note.

\section{Idea 1: is $q^\ast$ a period, a Nesterenko-type value, or a singular modulus?}

\textbf{Setup.} $q^\ast\in(0,1)$ is defined as the least positive root of
\[
S(q) = 1+\sum_{k\ge1}(-1)^k q^{2k(k-1)}\frac{[2q(1-q)]^k}{\prod_{j=1}^k(1-q^{2j-1})(1-q^{2j})},
\]
identified (paper2a, \texttt{beta2\_qniven\_ladder.py}) with a Hahn--Exton $q$-Bessel function
$J^{(3)}_{-1/2}(\cdot;q^2)$ up to reparametrization. Three concrete frameworks were considered
for whether $q^\ast$ itself, as opposed to the pair $(S,q^\ast)$, could be attacked
independently of a $\rho$-type exponent:

\begin{itemize}
\item \textbf{Nesterenko's theorem (1996).} Nesterenko's theorem gives algebraic independence
of $\pi$, $E_2(\tau)$/$E_4(\tau)$/$E_6(\tau)$-type quasimodular values at an algebraic $\tau$ in
the upper half plane, or equivalently of $q,\,\prod(1-q^n)$ combinations built from the
classical Eisenstein series/modular discriminant. It is a theorem about the transcendence
\emph{degree of a fixed triple of numbers built from $q$}, not a tool for showing a specific
$q$ satisfying an equation is itself transcendental --- it needs $q=e^{2\pi i\tau}$ with $\tau$
algebraic \emph{as an input}, and produces transcendence of \emph{values of Eisenstein series
at that q}, not information about which $q\in(0,1)$ solves an unrelated equation $S(q)=0$.
Applying it here would require going the other direction: showing $q^\ast=e^{2\pi i \tau^\ast}$
for some algebraic $\tau^\ast$ (which would make $q^\ast$ itself \emph{algebraic}, since
$e^{2\pi i\tau}$ for algebraic $\tau\ne0$ is transcendental by Gelfond--Schneider unless $\tau$
is rational, and $q^\ast\in(0,1)$ real forces $\tau^\ast$ purely imaginary, i.e.\
$q^\ast=e^{-2\pi t}$ for some real $t$). \textbf{HEURISTIC/diagnosis:} there is no evident reason
$q^\ast$ should be of the special form $e^{-2\pi t}$ with $t$ algebraic or even a period; nothing
in the defining series $S$ ties $q$ to an exponential-of-algebraic ansatz. Nesterenko's theorem
is the wrong tool: it is a statement about \emph{fixed, given} algebraic $\tau$, and $q^\ast$ is
not given as $e^{2\pi i\tau}$ for any distinguished $\tau$ --- it is defined as a zero of a
transcendental (in $q$) function. This route does not apply without additional structure that is
not present. \textbf{Conclusion: does not pan out, and I can state precisely why} (wrong
direction of implication; Nesterenko needs $\tau$ as the algebraic input, not as an unknown to be
solved for).

\item \textbf{Periods (Kontsevich--Zagier sense).} A period is (informally) the value of an
absolutely convergent integral of an algebraic function over a domain cut out by polynomial
inequalities with algebraic (in practice rational) coefficients. $q^\ast$ is defined as a root
of a power series equation, not as an integral. One can ask whether $S(q)=0$ can be rewritten as
a period-theoretic statement, e.g.\ via a contour-integral representation of the Hahn--Exton
$q$-Bessel function (such representations exist, e.g.\ Jackson-integral / $q$-Mellin--Barnes
forms with $q$ itself appearing as a parameter inside the integrand, not as the endpoint or the
object being solved for). \textbf{HEURISTIC/diagnosis:} even if $S(q)$ for \emph{fixed} $q$ were
exhibited as (a $q$-analogue of) a period-like integral, that would classify \emph{values of
$S$ at algebraic $q$} as periods, which is the opposite of what is needed here: we need
$q^\ast$, the location of a zero, to be shown transcendental, and ``the zero-locus of a period
function'' is not itself generally known to consist of periods or of any specific arithmetic
class --- zero sets of period-valued functions are exactly as hard as, e.g., zero sets of
$L$-functions (cf.\ the open-ness of transcendence of nontrivial zeros of $\zeta$-adjacent
objects). \textbf{Conclusion: the periods framework reclassifies the wrong object} (values of
$S$, not zeros of $S$); it does not bypass the original difficulty, it relocates it.

\item \textbf{Singular moduli / CM points.} If $q^\ast=e^{-\pi\sqrt N}$ (a nome at a CM point)
for some rational $N>0$, then $q^\ast$ would be transcendental by the Chudnovsky/Nesterenko-type
results on periods of CM elliptic curves (in fact by Schneider's classical theorem on singular
moduli, $j(\tau)$ algebraic $\Rightarrow$ strong constraints, and the associated nome
$q=e^{2\pi i\tau}$ for $\tau=i\sqrt N$ is transcendental by Gelfond--Schneider, since $\tau$ is
then a quadratic irrational hence algebraic and $\ne$ rational, so $q^\ast=e^{2\pi i\tau^\ast}$
would be transcendental \emph{immediately}, with no $\rho$-exponent needed at all). This is
the one genuinely promising bypass: if $q^\ast$ happened to equal $e^{-\pi\sqrt N}$ for a
positive rational $N$, transcendence of $q^\ast$ (hence of $\beta_2=q^{\ast-1/2}$) would follow
in one line from Gelfond--Schneider, with no Diophantine-approximation exponent problem at all.
\textbf{This was checked numerically below (Idea 2) and found NOT to hold} to 60-digit
precision for any small rational $N$ (nor did \texttt{mpmath.identify} find any relation of
$\tau^\ast:=-\ln(q^\ast)/\pi$ to a short list of standard constants). So the bypass is
structurally available in principle but numerically excluded for $q^\ast$ as it actually is.
\end{itemize}

\textbf{Net assessment of Idea 1:} none of the three frameworks apply to $q^\ast$ as defined.
Nesterenko's theorem needs $\tau$ as an algebraic input, not as something to be solved for, so
it is inapplicable in the direction needed. The periods framework reclassifies values of $S$,
not its zero-locus, and does not touch the actual difficulty. The singular-modulus/CM bypass
would work immediately \emph{if} $q^\ast$ were of that special exponential form, but the
numerics (Idea 2) say it is not, at least not with small $N$. \textbf{This is a genuine new
diagnosis} (not a reparametrization of the $\rho=1/2$ Diophantine-approximation route): the
obstruction is not merely ``$\rho$ is too small,'' it is that \emph{no known transcendence
machinery currently classifies numbers defined as zero-loci of $q$-Bessel-type functions in
their own base variable at all} --- not $q$-irrationality criteria (which need $\rho\ge1$), not
Nesterenko (needs $\tau$ given, not solved-for), not periods (classifies values, not zero
loci), and not Schneider-type CM results (needs the zero to happen to be a nome at a CM point,
which it numerically is not). This sharpens, rather than resolves, report\_G1.md's observation
that ``there is no base-variable theory for $q$-series that are neither modular nor Mahler'': the
present analysis identifies \emph{three specific named theories} that were checked and shown
inapplicable for three different, precise reasons, rather than leaving ``no theory applies'' as
an unexamined blanket statement.

\section{Idea 2: numerical recognition check on $q^\ast\approx0.4495$}

\textbf{VERIFIED (this session, mpmath, 60-digit working precision, using the pre-existing
2090-digit-certified value of $q^\ast$ from \texttt{private/ROOM/Beta2/g1\_certificate/bracket.txt}
truncated to 60 digits, not re-derived from the series in this session):}
\[
q^\ast = 0.449453630558948046125545825395706089225112808545993877680333\ldots,\qquad
\beta_2 = 1/\sqrt{q^\ast}.
\]
Sanity check performed: $1/\sqrt{q^\ast}-\beta_2 \approx 4.2\times10^{-60}$ (consistent to
full working precision).

\texttt{mpmath.identify} was run against the basis
$\{\pi,e,\sqrt2,\sqrt3,\sqrt5,\varphi,\log2,\log3,\gamma_{\mathrm{Euler}},\text{Catalan},\zeta(3)\}$
for both $q^\ast$ and $\beta_2$: \textbf{no relation found} (returned \texttt{None}) at 60-digit
precision for either constant.

\textbf{CM/nome check.} Set $\tau^\ast := -\ln(q^\ast)/\pi = 0.254559606032920872423\ldots$. If
$q^\ast=e^{-\pi\sqrt N}$ for rational $N$, then $\tau^{\ast2}=N$ would be rational; computed
$\tau^{\ast2}=0.0648005930236358846276\ldots$, which is not close to any small rational
$p/q$ with small $q$ (checked by continued-fraction expansion of $\tau^{\ast2}$; no small-denominator
convergent matches to 15+ digits, i.e.\ $\tau^{\ast2}$ does not look rational). \texttt{mpmath.identify}
on $\tau^\ast$ itself against the same constant basis: also \texttt{None}. As an additional sanity
check on the CM-point method itself (not on $q^\ast$), the true lemniscatic case $\tau=1$ was
computed via elliptic theta functions ($k(\tau)=(\theta_2/\theta_3)^2$ at nome $q=e^{-\pi}$),
correctly reproducing $k^2=1/2$ (the known lemniscatic value) --- confirming the numerical method
used is sound, and its negative finding for $q^\ast$ is not a tooling artifact.

\textbf{Conclusion of Idea 2:} $q^\ast$ does not match, to 60-digit precision, any of: (a) a
simple algebraic combination of $\{\pi,e,\sqrt2,\sqrt3,\sqrt5,\varphi,\log2,\log3,\gamma,G,\zeta(3)\}$
via \texttt{identify}'s search (bounded-degree/bounded-height rational combinations); (b) a nome
$e^{-\pi\sqrt N}$ at a small rational $N$ (the CM/singular-modulus bypass of Idea 1). Both checks
are negative, i.e.\ they rule out the two most obvious ``hidden structure'' hypotheses rather than
finding one. This is a real (if negative) result: it closes off the CM bypass identified in Idea
1 as numerically unavailable for $q^\ast$ as defined, so that particular one-line
Gelfond--Schneider argument cannot be used.

\section{Idea 3 (briefly considered, not pursued): Schneider--Lang}

Schneider--Lang requires two algebraically independent meromorphic functions of order $\le\rho$
satisfying a system such that some derivative/composition data is algebraic at sufficiently many
points, and concludes transcendence of a point where extra algebraic relations would otherwise
be forced. $S(q)$ is a single function of one base variable; no natural second, algebraically
independent function of finite order sharing algebraic-valued data with $S$ at a growing set of
points was found or constructed in this session. \textbf{HEURISTIC, brief:} Schneider--Lang's
machinery is built for exponential/elliptic-function pairs (e.g.\ $e^z,\wp(z)$) where algebraic
independence and finite order are both easy to certify; $S$'s $q$-series structure (order
considerations for a $q$-Bessel-type function as $q\to1^-$ vs.\ its behavior as a formal/base
variable) does not obviously fit this template, and constructing a genuine second function would
itself likely be as hard as the original problem. Not pursued further; flagged here only so the
route is documented as considered-and-set-aside rather than silently skipped.

\section{Summary verdict}

\begin{proposition}[status, mixed labels, this note]
\begin{enumerate}
\item[(a)] \textbf{HEURISTIC/diagnosis (new this session):} Nesterenko's theorem, the
Kontsevich--Zagier periods framework, and Schneider--Lang do not apply to $q^\ast$ as defined,
each for a distinct, statable reason (wrong direction of the algebraic-input hypothesis;
classifies values not zero-loci; no natural second function of matching order found).
\item[(b)] \textbf{VERIFIED (this session):} the one bypass that would work immediately if it
held --- $q^\ast$ being a CM nome $e^{-\pi\sqrt N}$, $N\in\mathbb Q_{>0}$ --- is numerically
excluded to 60-digit precision for small $N$; the numerical method was cross-validated on the
true lemniscatic case ($N=1$).
\item[(c)] \textbf{VERIFIED (this session):} no relation of $q^\ast$, $\beta_2$, or
$\tau^\ast=-\ln(q^\ast)/\pi$ to a standard basis of constants was found by \texttt{mpmath.identify}.
\item[(d)] \textbf{OPEN, and now more precisely scoped:} the obstruction to P6 is not solely the
$\rho=1/2$ deficit of \texttt{P6\_note.tex} --- it is the absence of \emph{any} applicable named
transcendence framework (four were checked here: Nesterenko, periods, singular moduli/CM,
Schneider--Lang; five Diophantine-approximation reparametrizations were checked in the prior
note) for numbers defined as zero-loci, in the base variable, of $q$-Bessel-type $q$-series. This
restates report\_G1.md's ``no base-variable theory'' diagnosis with four specific named
candidate theories now explicitly examined and excluded, rather than left as an unexamined gap.
\end{enumerate}
\end{proposition}

\textbf{Honest bottom line:} no proof was found in this session, and none of the ``genuinely
different'' frameworks solicited by the task (periods, Nesterenko, Schneider--Lang, CM/singular
moduli via numerical recognition) pan out for $q^\ast$/$\beta_2$ as they are actually defined.
The value of this note is diagnostic, not constructive: it (i) rules out the CM/nome bypass
numerically, closing off the one framework that would have given an immediate one-line proof
had it applied, and (ii) gives a precise, checkable reason each of the other three named
frameworks fails to apply, sharpening ``no base-variable theory exists'' from a general
observation into four specific closed doors.

\section{Weak points of this note}

\begin{itemize}
\item The CM/nome check (\S2) was only run against small rational $N$ and a short constant
basis for \texttt{identify}; it does not rule out $q^\ast$ being $e^{-\pi\sqrt N}$ for some
large or non-obvious rational $N$, nor a more exotic CM-type transcendental combination outside
the basis tried.
\item The claim that Nesterenko's theorem is inapplicable (\S1) rests on the standard statement
of the theorem as recalled from training and is flagged \textbf{UNVERIFIED-CITATION}: the exact
hypotheses of Nesterenko 1996 were not re-checked against a primary source in this session.
\item The Schneider--Lang discussion (\S3) is a brief plausibility argument, not an exhaustive
search for a suitable second function; a more determined search (e.g.\ trying the second
Hahn--Exton solution, or a Wronskian partner) was not attempted here for time/compute reasons.
\item As in the prior note, this is a diagnosis of why known machinery does not apply, not a
proof that no machinery ever could; ``genuinely new arithmetic input'' (per \S5 of the prior
note) is still what is required, and this note does not supply it.
\end{itemize}

\end{document}
